\documentclass[a4paper,titlepage,onecolumn,openright,twoside]{report}
\usepackage{amsmath}
\usepackage{graphicx} % required for pdf transparency
\usepackage{xcolor} % required for pdf transparency
\usepackage{amsthm}
\usepackage{amssymb}
\usepackage{booktabs}
\usepackage{tabularx}
\usepackage{subfigure}
\usepackage{mathrsfs}
\usepackage{textcomp}
\usepackage{siunitx}
\usepackage{fullpage}
\usepackage[multiple]{footmisc}
\usepackage{hyperref}
\usepackage{multirow}
\usepackage{empheq}
\usepackage{calc}
\usepackage[version=3]{mhchem}
\usepackage{tikz}
\usepackage{pgfplots}
\usepackage{xcolor}
\usepackage[section]{placeins}
\pgfplotsset{compat=newest}


%\hypersetup{
% colorlinks=false,
% linkbordercolor={1 1 1},
% citebordercolor={1 1 1},
% urlbordercolor={1 1 1} 
%}

\sisetup{ 
% load-configurations=abbreviations
% round-mode = places,
}%

% New definition of square root:
% it renames \sqrt as \oldsqrt
% This definition puts a little vertical guy at the end so it's more
% obvious where the square root actually ends.
\let\oldsqrt\sqrt
% it defines the new \sqrt in terms of the old one
\def\sqrt{\mathpalette\DHLhksqrt}
\def\DHLhksqrt#1#2{%
\setbox0=\hbox{$#1\oldsqrt{#2\,}$}\dimen0=\ht0
\advance\dimen0-0.2\ht0
\setbox2=\hbox{\vrule height\ht0 depth -\dimen0}%
{\box0\lower0.4pt\box2}}

% integrals with infinity bounds
\newcommand{\intinfty}{\int_{-\infty}^{\infty}}

% consistent formatting of object labels
\newcommand{\Figure}[1]{Figure \ref{#1}}
\newcommand{\Equation}[1]{Equation \ref{#1}}
\newcommand{\Table}[1]{Table \ref{#1}}
\newcommand{\Section}[1]{Section \ref{#1}}
\newcommand{\Chapter}[1]{Chapter \ref{#1}}
\newcommand{\Appendix}[1]{Appendix \ref{#1}}

% missing mathematical operators
\DeclareMathOperator{\sinc}{sinc}
\DeclareMathOperator{\sech}{sech}
\DeclareMathOperator{\sgn}{sgn}
\DeclareMathOperator{\erf}{erf}
\DeclareMathOperator{\inverf}{inverf}
\DeclareMathOperator{\arcsinh}{arcsinh}
\DeclareMathOperator{\arccosh}{arccosh}
\DeclareMathOperator{\arctanh}{arctanh}
%\DeclareMathOperator{\Re}{Re}
%\DeclareMathOperator{\Im}{Im}

% use roman type for natural base e and sqrt(-1)
\newcommand{\me}{{\mathrm{e}}}
\newcommand{\mi}{{\mathrm{i}}}

% roman type for the derivative, plus a space
\newcommand{\md}{\,\mathrm{d}}

% fourier transform and the reverse
\newcommand{\ff}[1]{{\mathscr{F}^{+}\bigl(#1\bigr)}}
\newcommand{\fr}[1]{{\mathscr{F}^{-}\bigl(#1\bigr)}}

% hilbert transform and the reverse
\newcommand{\hf}[1]{{\mathscr{H}^{+}\bigl(#1\bigr)}}
\newcommand{\hr}[1]{{\mathscr{H}^{-}\bigl(#1\bigr)}}

% custom lengths for figures

% width for side by side figures
\newlength{\twoupwidth}
\setlength{\twoupwidth}{7.5cm}

% width and height for default single figure
\newlength{\oneupwidth}
\setlength{\oneupwidth}{0.90\textwidth}
\newlength{\oneupheight}
\setlength{\oneupheight}{0.55623059\textwidth}

% custom colors for 2D plots - these are the same ones mathematica uses by
% default
\definecolor{colora}{RGB}{63,61,153}
\definecolor{colorb}{RGB}{153,61,113}
\definecolor{colorc}{RGB}{153,139,61}
\definecolor{colord}{RGB}{61,153,86}
\definecolor{colore}{RGB}{61,90,153}
\definecolor{colorf}{RGB}{153,61,144}
\definecolor{colorg}{RGB}{153,109,61}
\definecolor{colorh}{RGB}{67,153,61}
\definecolor{colori}{RGB}{61,121,153}
\definecolor{colorj}{RGB}{132,61,153}

% hilight text... a "work in progress" type thing
\tikzstyle{todobox} = [draw=red, fill=blue!10, very thick,
rectangle, rounded corners, inner sep=10pt, inner ysep=20pt]
\tikzstyle{todotitle} =[fill=red, text=white]

\newcommand{\todo}[1]{%
\begin{center}
\begin{tikzpicture}
\node [todobox] (box){%
 \begin{minipage}{0.9\textwidth}
 #1
 \end{minipage}
};
\node[todotitle, right=10pt] at (box.north west) {To Do:};
\end{tikzpicture}
\end{center}

}

\pgfplotsset{
 /pgfplots/colormap={jet}{rgb255(0cm)=(0,0,128) rgb255(1cm)=(0,0,255)
 rgb255(3cm)=(0,255,255) rgb255(5cm)=(255,255,0) rgb255(7cm)=(255,0,0)
 rgb255(8cm)=(128,0,0)}
}
\usepgfplotslibrary{units}
\usetikzlibrary{pgfplots.units} 


\pdfpageattr {/Group << /S /Transparency /I true /CS /DeviceRGB>>} 

\begin{document}

%%%%% TITLE PAGE %%%%%
\begin{titlepage}
\begin{center}
\hfill\\[4cm]
{ \Huge {\bfseries {Interference and Scattering in Surface Plasmon Resonance}} \par}
\vspace{3.0cm}
\begin{tabular}{lr}
\includegraphics[height=2cm,keepaspectratio]{logo/Logo_MPL_englisch_kompakt_cmyk_110915}
\hspace{1.0cm}
%
\includegraphics[height=2cm,keepaspectratio]{logo/Logo_IMPRS_4c_042012}
\end{tabular}
\vspace{3cm}
\\
{\Large Aaron Webster}\\
\vspace{1cm}
{\large A DISSERTATION}\\
\vspace{0.5cm}
Presented to the Max Planck Institute for the Science of Light\\
and the University of Erlangen-N\"urnberg\\
in partial fulfillment of the requirements\\
for the degree of\\
Doctor of Philosophy\\
\vspace{0.5cm}
\today
\end{center}
\end{titlepage}

%%%%% DECLARATION OF ORIGINALITY %%%%%
\chapter*{Declaration of Originality}
I affirm that the work presented in this dissertation is, to the best of my
knowledge, original and my own, except as acknowledged in the
text. \\
\hfill\\[1cm]
Signed\hspace{0.25cm}\makebox[5cm]{\hrulefill}\hspace{0.25cm}(Aaron Webster)
\hfill\\[1cm]
Date\hspace{0.51cm}\makebox[5cm]{\hrulefill}\hspace{0.25cm}
\vspace{2cm}

%%%%% ADVISOR, CO-ADVISOR %%%%%
\section*{Evaluation Committee}
\subsection*{First Advisor}
Dr. Frank Vollmer\\
Laboratory of Biophotonics and Biosensing\\
Max Planck Institute for the Science of Light\\
G\"unther-Scharowsky-Str.1 / Bau 24\\
91058 Erlangen
\subsection*{Second Advisor}
Prof. Dr. Ulf Peschel\\
Institute of Optics, Information and Photonics\\
University Erlangen-N\"urnberg\\
G\"unther-Scharowsky-Str.1 / Bau 24\\
91058 Erlangen

\newpage
\chapter*{Acknowledgements}
The author would like to extend the following thanks

\begin{itemize}
\item to Jiapeng Huang for his enthusiasm and hard work in carrying out
many tedious but necessary experiments.
\item to Yuqiang Wu for determining out the proper protocols regarding the use
of DNA origami.
\item and to Stephen Gregory and R.P. Schumann whose pioneering work
inspired this dissertation.
\end{itemize}

%%%%% TABLE OF CONTENTS &&&&&
\pagenumbering{roman}
\tableofcontents

%%%%% ABSTRACT %%%%%
\begin{abstract}
(abstract is written last)
\end{abstract}
% the first part is in roman numbering, switch now to arabic
\pagenumbering{arabic}

\chapter{Foundations}
\section{Historical Perspective}
The theoretical groundwork for the existence of surface plasmon polaritons
(SPPs) was first introduced by Richie in his seminal 1957 paper
\textit{Plasma Losses by Fast Electrons in Thin Films}
\cite{ritchie1957plasma}.  Like any scientific work, Richie's was
incremental and has its roots in earlier theoretical proposals by Pines and
Bohm~\cite{bohm1951collective}~\cite{pines1952collective}.  Ultimately this
research functioned to explain the phenomena of sharp and spectrally narrow
energy losses observed in diffraction gratings by Wood in 1902, known as
``Wood's anomaly''.

Optical excitation of surface plasmons was made accessible through
pioneering work in the late 1960's by Kretschmann~\cite{kretschmann1968},
Raether~\cite{raether1965springer} and Otto~\cite{otto1968excitation}.
These experiments used the principle of attenuated total reflection (ATR)
to excite surface plasmons evanescently, using a prism to match their
resonance condition.  A great deal of understanding on the topic of surface
plasmons took place in the subsequent decade, such as an improved
theoretical understanding based on Fresnel relations
\cite{chen1976excitation} and descriptions of of conically scattered light
in the presence of surface roughness~\cite{simon1976directional}.  A
concise overview of this research can be found in~\cite{raether1997surface}.

The introduction of SPR as a biosensing platform began in the early 1980's
with work by Liedberg, Nylander and Lundstrom~\cite{liedberg1983surface}
who described the extraordinary sensitivity of the surface plasmon
resonance condition to perturbations in the refractive index of the medium
on one side of the film.  The subsequent commercialization of SPR
biosensors has largely been influenced by these authors and Pharmacia
Biosensor AB (now Biacore)~\cite{liedberg1995biosensing}.

The commercial success of biosensors based on surface plasmon resonance
seems to have brought about a knowledge gap between the biosensing
community and their more theoretical predecessors from whom the field owes
its genesis.  This is to say that the scope of SPR biosensing experiments
is disproportionately narrower than the breadth of phenomena discovered
since Richie.  As an example particular to this dissertation, in 2005 and
2007, two papers~\cite{andaloro2005optical}~\cite{simon2007observation}
based on theoretical work by Chuang~\cite{chuang1986lateral} and Chen
\cite{chen1976excitation} reported a curious interference pattern occurring
in the specularly reflected light for certain (among them,
Kretschmann-Raether type) systems illuminated with focused Gaussian beams.
This was also independently reported a year later in
\cite{schumann2008near}.  Interestingly, observation of this interference
required nothing more than the addition of a lens pair to a fairly
ubiquitous optical setup, but it somehow escaped attention during earlier
research.  

Many of the topics here are inspired by work done at the University of
Oregon in the labratory of Stephen Gregory, summarized in a 2009 thesis
{\it Surface Plasmon Random Scattering and Related Phenomena}
\cite{schumann2009surface} by R.P. Schumann.  Here is described experiments
performed on thin metal films in a Kretschmann-Raether configuration with
the addition of a scanning apertureless near-field probe.  This probe (a
sharp tungsten tip) is able to elastically scatter SPPs in a way analogous
to surface roughness, but in this case its location and interaction can be
precisely controlled.  

\section{Surface Plasmon Polaritons}
The physics of surface plasmon polaritions has been the the topic of many
excellent publications. I will therefore present a minimalist theoretical
description of SPPs and relevant mathematics for predicting their behavior
in multilayer systems.
\section{The Kretschmann Configuration}
\section{Conically Scattered Light}

\chapter{Interference}
Surface plasmon resonance is in of itself an interference phenomena: the
dark band in the specular direction is due to interference between
specularly reflected light and the re-radiated SPPs which are likewise
antiphase at the SPR resonance angle.  There is, however, a more rich 
interference structure which can observed when using a focussed beam whose
spot size is smaller than the SPP propagation distance. 

\section{The Fourier Optics Perspective}
\section{Bulk Sensitivity}
\subsection{Angular Interogation}
\subsection{Intensity Interrogation}
\subsection{Minimum Resolvable Shift}

\chapter{Multiple Scattering of SPPs}
\section{Speckle}
\section{Optical Vorticies}
% pseudophase algorithm

\chapter{The Moving Scatterer Problem}
\section{Weirdospace}
\section{Reverse Mapping}

%\import{wiggles/}{fig2}

%\begin{figure}
%\centering
%\includegraphics{tmp/out.pdf}
%\caption{How do the font sizes look?}
%\end{figure}

\bibliographystyle{plain}
\bibliography{bibliography}

\appendix
\chapter{Protocols}
\section{Sputtering}
\section{Functionalizing Nanoparticles}
\section{Microfluidics}

\chapter{Reference Data}
Several datasets were produced during the course of this dissertation which
might be useful to the reader.  They are presented here.
\section{Physical Properties of SPP Propagation}
%\import{dispersion/}{sptable_632.tex}

\section{The Three Layer Fresnel System}
%% you will have to fix the path in the following file
%\import{fresnel_plots/}{allthreelayerfresnel}

\section{Complex Permittivities and Permiabilities of Metals}
%\import{eps_plots/}{allmetals}


\section{Abbreviations}
\begin{table}
 \begin{tabular}
  SPP & surface plasmon polariton
  LRSPP & long range surface plasmon polariton
 \end{tabular}
\end{table}

\chapter{Miscellaneous}
\section{Colophon}
This document was typeset using \pdfTeX {\verb 3.1415926-2.5-1.40.14}. All
plots were rendered using pgfplots~\cite{feuersangerpgfplots} with styles influenced by
Eduard Tufte~\cite{tufte1983visual}~\cite{tufte1983visual}. Colors were chosen
from datasets researched by Jan Brewer~\cite{harrower2003colorbrewer}.
Schematics and other line drawings were prepared using Inkscape. Generation
of text was done under Arch Linux using \texttt{vim}.  Some tables were
pre-processed using LyX.

The Monte Carlo simulations were written in \texttt{c} with pararallelism
accomplished using \texttt{mpi}. The vortex tracking and image processing
algorithms were written in \texttt{octave}.  Weirdospace images were
processed using python and scipy/numpy.  Parallel execution was done at at
the FAU RRZE, Erlangen, Germany.

\end{document}
